O \textit{Gradient Boosting} é um método de aprendizado de máquina baseado em conjuntos de árvores de decisão que
constrói modelos preditivos sequencialmente. Esta técnica combina modelos fracos (tipicamente árvores rasas) para formar
um preditor robusto, onde cada novo modelo corrige os erros residuais do conjunto anterior \cite{friedman}. O processo
inicia-se com um modelo constante $F_0(x) = min_{\gamma} \sum_{i=1}^n \mathcal{L}(y_i, \gamma)$, geralmente a média
dos valores-alvo. 

Em cada iteração $m$, o algoritmo calcula o gradiente negativo da função de perda $\mathcal{L}$ em relação às predições
atuais: $r_{im} = -\left[ \frac{\partial \mathcal{L}(y_i, F_{m-1}(x_i))}{\partial F_{m-1}(x_i)} \right]$.
Subsequentemente, ajusta-se uma árvore de decisão $h_m(x)$ a esses resíduos e atualiza o modelo conforme $F_m(x) =
F_{m-1}(x) + \gamma_m h_m(x)$, onde $\gamma_m$ denota a taxa de aprendizado. 

Entre suas vantagens destacam-se a alta precisão preditiva, flexibilidade para dados heterogêneos, suporte a funções de
perda personalizadas e aplicabilidade em regressão e classificação. Contudo, apresenta sensibilidade a
\textit{overfitting} (controlável via \textit{early stopping}), demanda ajuste rigoroso de hiperparâmetros como
profundidade das árvores, taxa de aprendizado e número de iteradores, além de custo computacional elevado.